\section{回到长安，不等于收回军镇}
\label{sec:middle-tang-reconstruction}

779年德宗即位时，唐廷已经能够在长安重新发布诏令，却未能把河北军镇变回由京师任免、供给的普通州县。次年两税法以夏、秋征收取代难以贯彻的租庸调，意在把战乱后散失的户籍、土地和收入重新接入财政；李惟岳等在781年拒绝朝廷限制父死子继，随即使这种重建的限度显露出来。税法不是先有一个完整国库再向地方征取的技术调整，而是中央必须同拥有兵员、仓储和地方关系的军镇重新议价的条件。%
\footnote{德宗即位、两税法与河北军镇反叛见新增账本 LT-779-019、LT-780-020、LT-781-021；可参《旧唐书·德宗本纪》《旧唐书·杨炎传》《新唐书·食货志五》及《通典》。诏令说明征收原则，不能证明各道税率、户等和留截完全一致。}

783年泾原军哗变、朱泚据长安，德宗出奔奉天；李希烈又在蔡州称楚帝。次年《奉天罪己诏》以赦宥和调整政策争取回旋，785年朝廷得以复归京师。这里应当把“罪己”读作危机中的政治文书，而不只是帝王德行的表白：它需要让军镇、官员和地方势力相信，重新接上朝廷比继续战争更有利。平凉会盟中唐使遭吐蕃伏击，以及790年安西、北庭联络的终结，又显示中原的缓和并未消除西北的军事与交通代价。%
\footnote{泾原兵变、奉天诏、复都、平凉与安西北庭终结见 LT-783-022 至 LT-790-027，依据《旧唐书》德宗纪、朱泚传、吐蕃传及《资治通鉴》。会盟责任、赦令实际效果和西域联络的具体断绝过程，在汉文叙事与边疆材料中均有缺环。}

805年永贞改革试图罢宫市、限制宦官，旋因顺宗让位而夭折；这并不证明改革集团只是一段宫廷插曲。它说明市场征取、内廷权力和皇位继承已经彼此相连。宪宗随后在蜀地平刘辟、在淮西围攻吴元济，817年李愬入蔡州，是中央重新取得局部军事主动的结果；但朝廷仍需依靠诸道军、转运和地方供给，不能把一场夜袭写成军镇问题已被永久解决。%
\footnote{永贞改革、刘辟、淮西战争及蔡州被平见 LT-805-029、LT-805-030、LT-807-031、LT-813-033 至 LT-817-036，依据《旧唐书》顺宗、宪宗纪及王叔文、刘辟、吴元济、李愬诸传，并参《资治通鉴》。传记的忠奸格局和夜袭细节不等于可完整复原的军粮、兵数与地方反应。}

819年迎佛骨及韩愈被贬、821年科举争议与唐蕃会盟、828--829年沧景之乱，分别发生在宗教、取士、外交和军镇中，却都指向同一事实：朝廷的秩序不是一条由长安单向传出的命令。佛教物资和声望、考试中的门第与人情、碑约中的边界语言以及沧州军人的继承选择，都是不同层次的政治资源。碑刻、文集和正史能够保留这些争论的声音，也大多保留精英的声音；普通纳税者、军户与寺院依附者如何经历这番重建，仍缺少连续记录。%
\footnote{佛骨、科举争议、长庆会盟及沧景之乱见 LT-819-037、LT-821-039、LT-821-040、LT-828-042、LT-829-043。核心材料为《旧唐书》《新唐书》本纪、食货志与吐蕃传、《唐会要》及会盟碑；碑文和科场案的立场都须与其政治用途一并阅读。}
